\documentclass[margin,line]{res}

% File produced by Jeremy West
% This file may be distributed and/or modified:
%   1. under the LaTeX Project Public License and/or
%   2. under the GNU Public License.

% NOTE: you also need the (included) res.cls class file to compile this document
% Either place the .cls file in the same directory (folder) as this .tex file, or install it globally for your LaTeX distribution (search online for how to do so).

\usepackage[dvipsnames]{xcolor}
\usepackage[export]{adjustbox}
\usepackage{graphicx}
\usepackage{times}
\usepackage{setspace}                                    % Allows for custom margins, etc.
\usepackage{fullpage}                                    % Use the full page
\usepackage[latin9]{inputenc}                            % Font definition and input type
\usepackage[T1]{fontenc}                                 % Font output type
\usepackage{mmap}                                        % Allows glyphs (e.g. ff) to copy properly (ASCII)
\usepackage{textcomp}                                    % Supports many symbols such as copyright
\usepackage{color}                                       % Allow for colored text, etc.
\usepackage[]{hyperref}                                  % Allow hyperlinks (internal and external)
\hypersetup{                                             % Custom hyperlink settings
    pdffitwindow=true,                                  % Window fit to page when opened
    pdfstartview={XYZ null null 1.00},                   % Fits the zoom of the page to 100%
    pdfnewwindow=true,                                   % Links in new window
    colorlinks=true,                                     % false: boxed links; true: colored links
    linkcolor=black,                                     % Color of internal links (black is necessary for printing quality)
    citecolor=black,                                     % Color of links to bibliography
    urlcolor=[rgb]{0.31,0.16,0.52},                                      % Color of external links
    pdfauthor = {Modibo Camara},
    pdfkeywords = {Economics, Northwestern},
    pdftitle = {CV, Modibo Camara},
    pdfsubject = {Curriculum Vitae},
    pdfpagemode = UseNone}

\oddsidemargin -.5in
\evensidemargin -.5in
\textwidth=6.0in
\itemsep=0in
\parsep=0in
% if using pdflatex:
\setlength{\pdfpagewidth}{\paperwidth}
\setlength{\pdfpageheight}{\paperheight}

	\addtolength{\topmargin}{-.3in}
	\addtolength{\textheight}{0.6in}

\newenvironment{list1}{
  \begin{list}{\ding{113}}{%
      \setlength{\itemsep}{.025in}
      \setlength{\parsep}{0in} \setlength{\parskip}{0in}
      \setlength{\topsep}{0in} \setlength{\partopsep}{0in}
      \setlength{\leftmargin}{0.17in}}}{\end{list}}
\newenvironment{list2}{
  \begin{list}{$\bullet$}{%
      \setlength{\itemsep}{0in}
      \setlength{\parsep}{0in} \setlength{\parskip}{0in}
      \setlength{\topsep}{0in} \setlength{\partopsep}{0in}
      \setlength{\leftmargin}{0.2in}}}{\end{list}}

\usepackage[nodayofweek]{datetime}                                    % Custom date format for date field
\newdateformat{mydate}{\monthname[\THEMONTH] \THEYEAR}   % Defining month year date format

\usepackage{fancyhdr}                                    % Used for custom page headers
\pagestyle{fancy}
\fancyhf{}
\renewcommand{\headrulewidth}{0.5pt}
\rhead{\footnotesize Modibo Camara \thepage} %header at the right
\headsep = 0.5cm
% FIRST PAGE ONLY (redefine the plain pagestyle)
\fancypagestyle{plain}{
\fancyhf{}
\renewcommand{\headrulewidth}{0pt}
\headsep = 0.0cm
%\rhead{\footnotesize Last updated \today}
\rhead{\footnotesize \includegraphics[scale=.11]{Shorthand-Vertical-Purple}}
}


\begin{document}
\pretolerance=10000

~\\
~\\
\name{  {\LARGE  Modibo Camara} \vspace*{.1in}}

\begin{resume}
\thispagestyle{plain} % to use first page footer

\begin{table}[h]
\footnotesize
\begin{tabular}{@{}p{0.20in}p{1.5in}p{1.4in}p{.7in}p{1.5in}}
& Placement Director:	& Professor Alessandro Pavan & 847-491-8266  & \href{mailto: alepavan@northwestern.edu}{alepavan@northwestern.edu}\\
& Placement Administrator:	& Alison Stoute &	847-491-5694  & \href{mailto: econjobmarket@northwestern.edu}{econjobmarket@northwestern.edu}\\
 \end{tabular}
\end{table}

\section{ Contact Information}
\vspace{.05in}
\begin{tabular}{@{}p{0.20in}p{2.75in}p{2.75in}}
 & Department of Economics           &Mobile: 202-415-2751 \\
 & Northwestern University   & \href{mailto: mcamara@u.northwestern.edu}{mcamara@u.northwestern.edu} \\
 & 2211 Campus Drive & \href{https://sites.northwestern.edu/mkc6866/}{https://sites.northwestern.edu/mkc6866/}\\
 & Evanston, IL 60208  & Citizenship: United States, Germany
\end{tabular}

\section{ Fields}
\begin{list1}
\item[] Research: Microeconomic Theory, Economics \& Computation
\item[] Teaching: Microeconomics, Econometrics
\vspace*{.05in}
\end{list1}

\section{ Education}\begin{tiny}

\end{tiny}
\begin{list1}
\item[] Ph.D., Economics, Northwestern University\hfill (anticipated) 2022
\item[] \begin{list2}
		\item[] Dissertation: Complexities in Economic Theory
		\item[] Committee: Eddie Dekel (Chair), Jason Hartline (Chair), Marciano Siniscalchi, Jeffrey Ely
	\end{list2}
\item[] M.A., Economics, Northwestern University\hfill 2018
\item[] B.A., Mathematics \& Economics, University of Pennsylvania\hfill 2016
\end{list1}

\section{Fellowships \& Awards}
\begin{list1}
\item[] Distinguished TA Award\hfill 2021-22
\item[] Dissertation University Fellowship, Northwestern University\hfill 2021-22
\end{list1}

\section{Teaching Experience}
\begin{list1}
\item[] Teaching Assistant, Northwestern University\hfill 2017-2021
\item[] \begin{list2}
		\item[] Econometrics (graduate), Econometrics (undergrad), Microeconomics (undergrad)
	\end{list2}
\end{list1}

\section{Research Experience}
\begin{list1}
\item[] Research Intern, Nicole Immorlica, Microsoft Research\hfill 2020
\item[] Research Assistant, Eddie Dekel, Northwestern University\hfill 2019
\item[] Project Intern, Yian Liu, Federal Reserve Board of Governors\hfill 2015
\item[] Intern, Esen Onur, Commodity Futures Trading Commission\hfill 2014
\item[] Research Assistant, Susan Wachter, University of Pennsylvania\hfill 2013-14
%\vspace*{.05in}
\end{list1}

\section{ Conferences}
\begin{list1}
\item[] European Winter Meeting of the Econometric Society\hfill 2021
\item[] Cornell ORIE (invited)\hfill 2021
\item[] Economics and Computation (invited)\hfill 2021
\item[] Transatlantic Theory Workshop\hfill 2021
\item[] Risk, Uncertanty, and Decision (RUD)\hfill 2021
\item[] Decision: Theory, Experiments, and Applications (D-TEA)\hfill 2021
\item[] North American Summer Meeting of the Econometric Society \hfill 2021
\item[] Africa Meeting of the Econometric Society \hfill 2021
\item[] European Summer Meeting of the Econometric Society \hfill 2021
\item[] World Congress of the Game Theory Society (GAMES)\hfill 2021
\item[] Foundations of Computer Science (FOCS)\hfill 2020
\item[] European Winter Meeting of the Econometric Society\hfill 2020
\item[] Young Economists Symposium (YES)\hfill 2020
\item[] Economics Graduate Student Conference (EGSC)\hfill 2020
\item[] Midwest Theory Day\hfill 2018
\end{list1}

\section{Refereeing}
\begin{list1}
\item[] American Economic Review, Journal of Mathematical Economics, The Web Conference
\end{list1}

\section{Job Market Paper}
\begin{list1}
\item[] \textbf{``Computationally Tractable Choice''}
\vspace*{.05in}
	\begin{list2}
		\item[] \textit{Abstract:} 
		I incorporate computational constraints into decision theory in order to capture how cognitive limitations affect behavior.
		I impose an axiom of computational tractability that rules out behaviors that are thought to be fundamentally hard. 
		I use this framework to better understand common behavioral heuristics: if choices are tractable and consistent with the expected utility axioms, then they are observationally equivalent to forms of choice bracketing.
		Then I show that a computationally-constrained decisionmaker can be objectively better off if she is willing to use heuristics that would not appear rational to an outside observer.
	\end{list2}
\end{list1}

\section{Other Papers}
\begin{list1}
\item[] \textbf{``Mechanisms for a No-Regret Agent: Beyond the Common Prior''} (with  J. Hartline, A. Johnsen)
\vspace*{.05in}
	\begin{list2}
		\item[] \textit{Proceedings of the IEEE Symposium on Foundations of Computer Science (FOCS 2020)}
		\vspace*{.05in}
		\item[] \textit{Abstract:} 
		We study incomplete information games where a policymaker commits to a policy, an agent responds, and payoffs depend on a state of nature.
		To solve these games, it is typical to (i) make strong prior knowledge assumptions or (ii) make no prior knowledge assumptions and optimize against the worst case.
		We show that, if the designer and agent interact repeatedly over time, we can approximate the superior performance of (i) while preserving the robustness of (ii).
		The state is revealed after every period, but we make no other assumptions on the state-generating process.
		We develop calibrated policies that adapt to historical data over time, assuming the agent does the same, even if the data is highly non-stationary.
		This requires new behavioral assumptions that build on prior work on learning in games and capture ideas like ``rationality'' and ``unpredictability'' in a fully ex post sense.
		
%		We study repeated games of incomplete information between a policymaker with commitment power and a single agent.
%		We propose policies that adapt to historical data over time, assuming the agent does the same, without making any assumptions about the data-generating process.
%		This requires new behavioral assumptions that capture ideas like ``rationality'' and ``unpredictability'' in a fully ex post sense.
%		Our policies are competitive with optimal static policies that rely on much stronger assumptions, like common prior beliefs.
		
	\end{list2}
\end{list1}
\vspace*{.1in}
\begin{list1}
	\item[] \textbf{``Mechanism Design with a Common Dataset''}
\vspace*{.05in}
	\begin{list2}
		\item[] \textit{Abstract:} 
		I propose a new approach to mechanism design: rather than assume a common prior belief, assume access to a common dataset.
		I restrict attention to incomplete information games where a designer commits to a policy and a single agent responds.
		I propose a penalized policy that performs well under weak assumptions on how the agent learns from data.
		Policies that are too complex, in a precise sense, are penalized because they lead to unpredictable responses by the agent.
		This approach leads to new insights in models of vaccine distribution, prescription drug approval, performance pay, and product bundling.
	\end{list2}
\end{list1}

\section{Languages}
\begin{list1}
\item[] English (fluent), German (native), Spanish (intermediate)
\end{list1}

\section{ References}
\vspace{.05in}
%Pick from:
\begin{tabular}{@{}p{0.20in}p{2.75in}p{2.75in}}
 & Professor Eddie Dekel           & Professor Jason Hartline \\
 & Department of Economics   & Department of Computer Science \\
 & Northwestern University   & Northwestern University \\
 & 2211 Campus Drive  & 2233 Tech Drive  \\
 & Evanston, IL 60208  & Evanston, IL 60208\\
 & 847.491.4414  & 847.467.0280\\
 &  \href{mailto: dekel@northwestern.edu}{dekel@northwestern.edu} &  \href{mailto: hartline@northwestern.edu}{hartline@northwestern.edu} \\
\end{tabular}
\vspace{.1in} ~\\
\begin{tabular}{@{}p{0.20in}p{2.75in}p{2.75in}}
	& Professor Marciano Siniscalchi           & Professor Jeffrey Ely \\
	& Department of Economics   & Department of Economics \\
	& Northwestern University   & Northwestern University \\
	& 2211 Campus Drive  & 2211 Campus Drive  \\
	& Evanston, IL 60208  & Evanston, IL 60208\\
	& 847.491.5398  & 847.491.8208\\
	&  \href{mailto: marciano@northwestern.edu}{marciano@northwestern.edu} &  \href{mailto: jeffely@northwestern.edu}{jeffely@northwestern.edu} \\
\end{tabular}
%\begin{tabular}{@{}p{0.20in}p{2.75in}p{2.75in}}
% & Professor XX            \\
% & Department of Economics   \\
% & Northwestern University    \\
% & 2211 Campus Drive    \\
% & Evanston, IL 60208  \\
% & 847.491.xxxx  \\
% &  \href{mailto: xx@northwestern.edu}{xx@northwestern.edu}  \\
%\end{tabular}



\end{resume}
\end{document}
